\documentclass[11pt]{article}
\usepackage[localtheorems]{raeez-math-template}

\newtheorem{theorem}{Theorem}[section]
\newtheorem{proposition}[theorem]{Proposition}
\newtheorem{lemma}[theorem]{Lemma}
\newtheorem{definition}[theorem]{Definition}
\newtheorem{conjecture}[theorem]{Conjecture}
\newtheorem{remark}[theorem]{Remark}

\providecommand{\C}{\mathbb{C}}
\providecommand{\Z}{\mathbb{Z}}
\providecommand{\Q}{\mathbb{Q}}
\providecommand{\CoHA}{\mathrm{CoHA}}
\providecommand{\Rep}{\mathrm{Rep}}
\providecommand{\kcat}{\kappa_{\mathrm{cat}}}
\providecommand{\kch}{\kappa_{\mathrm{ch}}}
\providecommand{\kBKM}{\kappa_{\mathrm{BKM}}}
\providecommand{\kfib}{\kappa_{\mathrm{fiber}}}

\title{The \(E_8\times E_8\) Yangian and the Super-Yangian Target}
\author{Raeez Lorgat}
\date{2026-04-22}

\begin{document}
\maketitle

\begin{abstract}
The heterotic \(E_8\times E_8\) current algebra, the K3 Mukai lattice,
and the \(\Delta_5\) Borcherds denominator determine three related
objects with different proofs.  The Drinfeld Yangian
\(Y_\hbar(\mathfrak e_8\oplus\mathfrak e_8)\) is standard.  The
degree-\((24,24)\) rational function is the Maulik--Okounkov geometric
\(R\)-matrix on the K3 Mukai Fock representation.  The super-Yangian
\(Y(\mathfrak{gl}(4|20))\) and its embedding into the completed
Hall--Drinfeld/Borcherds double are conjectural.
\end{abstract}

\tableofcontents

\section{The Standard Yangian}
\label{sec:e8-standard-yangian}

Let \(\mathfrak g=\mathfrak e_8\oplus\mathfrak e_8\).  The Drinfeld
Yangian \(Y_\hbar(\mathfrak g)\) is generated by
\[
  x^{\pm}_{i,r},\qquad h_{i,r},\qquad
  i\in\{1,\ldots,16\},\quad r\geq0,
\]
with Drinfeld's current relations and the simply-laced Serre families.
The direct-sum identity is
\[
  Y_\hbar(\mathfrak e_8\oplus\mathfrak e_8)
  \cong
  Y_\hbar(\mathfrak e_8)\otimes_{\C[\hbar]}
  Y_\hbar(\mathfrak e_8).
\]

\begin{proposition}[Dynkin edge count]
\label{prop:e8-edge-count}
The finite Dynkin edge count for
\(\mathfrak e_8\oplus\mathfrak e_8\) is \(14\), with \(7\) edges in
each \(E_8\) factor.
\end{proposition}

\begin{proof}
The \(E_8\) Dynkin diagram is a tree on \(8\) vertices, hence has
\(7\) edges.  The direct sum has two disconnected \(E_8\) diagrams, so
the edge count is \(7+7=14\).  Drinfeld's Serre families are attached to
these edge types, together with the sign and loop-mode data of the
current presentation.
\end{proof}

\section{Intrinsic and Geometric Rational Functions}
\label{sec:intrinsic-geometric}

The intrinsic Yangian current relation has one linear spectral factor
per Cartan entry:
\[
  g_{ij}(u)
  =
  \frac{u+\hbar a_{ij}/2}{u-\hbar a_{ij}/2}.
\]
For \(\mathfrak e_8\oplus\mathfrak e_8\) the finite Cartan rank is
\(16\).  The rank-\(24\) rational function arises after passing to the
K3 Mukai representation.

\begin{proposition}[Mukai \(R\)-matrix degree]
\label{prop:mukai-rmatrix-degree}
On the K3 Nakajima--Mukai Fock representation, the
Maulik--Okounkov stable-envelope \(R\)-matrix has a product
normalisation with one spectral factor for each basis class in
\[
  \widetilde H(K3,\Q)
  =
  H^0(K3,\Q)\oplus H^2(K3,\Q)\oplus H^4(K3,\Q),
  \qquad
  \dim \widetilde H(K3,\Q)=24 .
\]
In this representation its numerator and denominator have degree
\((24,24)\).
\end{proposition}

\begin{proof}
The Nakajima construction identifies the K3 Hilbert-scheme Fock space
with the symmetric algebra generated by
\(\widetilde H(K3,\Q)\otimes t^{-1}\C[t^{-1}]\).  Stable envelopes give
wall \(R\)-matrices whose diagonal factors are indexed by the Heisenberg
oscillator labels.  The primitive oscillator labels form a basis of the
rank-\(24\) Mukai lattice, so the product normalisation carries \(24\)
linear numerator factors and \(24\) linear denominator factors.
\end{proof}

\section{Central Charge}
\label{sec:central-charge}

The level-one affine current algebra has
\[
  c(\widehat{\mathfrak e}_8)_1
  =
  \frac{\dim\mathfrak e_8}{1+h^\vee_{E_8}}
  =
  \frac{248}{31}
  =
  8.
\]
Thus the two heterotic gauge factors contribute \(8+8=16\).  The
heterotic left-moving matter theory adds eight transverse bosons:
\[
  c_L^{\mathrm{matter}}=8+8+8=24 .
\]
This is the worldsheet matter central charge.  It is separate from the
intrinsic Yangian presentation.

\section{The \(4|20\) Mukai Target}
\label{sec:4-20-target}

The real Mukai lattice has signature \((4,20)\):
\[
  \widetilde H(K3,\Z)\cong II_{4,20}.
\]
The positive four-plane is spanned by the real and imaginary parts of
the holomorphic \(2\)-form together with the positive plane in
\(H^0\oplus H^4\).  Its orthogonal complement has rank \(20\).  This
signature gives the natural super-dimension for the proposed
super-Yangian target:
\[
  \mathfrak{gl}(4|20).
\]

\begin{definition}[Super-Yangian target]
\label{def:super-yangian-target}
Let \(V=V_{\bar0}\oplus V_{\bar1}\) with
\(\dim V_{\bar0}=4\) and \(\dim V_{\bar1}=20\).  The RTT
super-Yangian \(Y(\mathfrak{gl}(4|20))\) is generated by the coefficients
of a \(24\times24\) monodromy matrix \(T(u)\), with graded RTT relation
\[
  R(u-v)T_1(u)T_2(v)=T_2(v)T_1(u)R(u-v),
  \qquad
  R(u)=1+\frac{\hbar P_{\mathrm{super}}}{u}.
\]
\end{definition}

\begin{conjecture}[Hall--Borcherds super-Yangian bridge]
\label{conj:super-yangian-bridge}
There is a completed Hopf-superalgebra morphism
\[
  Y(\mathfrak{gl}(4|20))
  \longrightarrow
  \mathbf D_\hbar\bigl(\CoHA^+_{K3\times E}\bigr)
\]
whose image is the Mukai-finite sector of the compact
Hall--Drinfeld/Borcherds double.  Under finite BKM recognition, this
image maps to the corresponding finite sector of
\(\mathfrak g_{\Delta_5}\).
\end{conjecture}

The conjecture requires the compact positive half, the Serre-dual
negative half, the Hopf pairing, completion, PBW/radical quotient, and
Borcherds recognition.  The \(E_8(-1)^2\) lattice supplies a heterotic
rank-\(16\) sublattice inside the \(20\)-dimensional negative sector.

\section{Six-Dimensional Boundary}
\label{sec:six-dimensional-boundary}

A six-dimensional hCS construction with heterotic boundary has two
layers:
\[
  \widehat{\mathfrak e}_8{}_1\oplus\widehat{\mathfrak e}_8{}_1
  \quad\leadsto\quad
  Y_\hbar(\mathfrak e_8\oplus\mathfrak e_8),
\]
and the conjectural Mukai-enhanced layer
\[
  \widetilde H(K3)\quad\leadsto\quad Y(\mathfrak{gl}(4|20)).
\]
The compact \(K3\times E\) comparison also needs the usual hCS package:
orientation and determinant-line framing, regulator, completed local
functionals, QME solution, quartic anomaly trivialisation, and
factorisation or TCFT sewing.  The \(d=3\) CY-to-chiral output is
\(E_1\)-chiral.  Braided \(E_2\) structure belongs to a centre, double,
or module layer.

\section{K3 \(\times\) E Scalars}
\label{sec:k3xe-scalars}

For \(X=K3\times E\) the relevant scalars are
\[
  \kcat(X)=0,\qquad
  \kch^{\mathrm{compact}}(X)=0,\qquad
  \kch^{\mathrm{Heis}}(X)=3,\qquad
  \kBKM(\mathfrak g_{\Delta_5})=5,\qquad
  \kfib(K3)=24.
\]
They arise from the total-space Euler characteristic, compact Hodge
supertrace, Heisenberg-Mukai specialisation, Borcherds weight, and
Mukai rank respectively.  The heterotic \(E_8\times E_8\) current
algebra contributes \(16\) gauge central-charge units inside a separate
worldsheet computation.

\section{Obligations}
\label{sec:obligations}

The proved statements are the standard Yangian
\(Y_\hbar(\mathfrak e_8\oplus\mathfrak e_8)\), the \(E_8\) edge count,
the Sugawara central charge \(8+8\), and the rank-\(24\) Mukai Fock
input to the geometric \(R\)-matrix.  The super-Yangian bridge requires:
\begin{enumerate}[label=\textup{(\arabic*)}]
\item a compact \(K3\times E\) Hall positive half;
\item a Serre-dual negative half and Hopf pairing;
\item a completed Hall--Drinfeld double;
\item a PBW theorem and radical quotient;
\item a morphism from the RTT super-Yangian;
\item finite Borcherds recognition against \(\mathfrak g_{\Delta_5}\);
\item compatibility with the \(E_1\) chiral boundary and its centre or
double.
\end{enumerate}

\end{document}
